\documentclass[11pt]{article}
\usepackage{fontspec}
\setmainfont{DejaVu Serif}
\usepackage[margin=1.8cm]{geometry}
\usepackage{longtable}
\usepackage{array}
\usepackage{xcolor}
\usepackage{titlesec}
\pagestyle{plain}

\newcolumntype{L}[1]{>{\raggedright\arraybackslash}p{#1}}

\setlength{\tabcolsep}{5pt}
\renewcommand{\arraystretch}{1.25}
\setlength{\parindent}{0pt}
\setlength{\parskip}{5pt}

\titleformat{\section}{\Large\bfseries}{\thesection.}{0.5em}{}
\titleformat{\subsection}{\large\bfseries}{\thesubsection.}{0.5em}{}

\newcommand{\qa}[2]{%
\par\vspace{7pt}%
{\bfseries Вопрос: #1}\par\vspace{2pt}
{Ответ: #2}\par\vspace{4pt}}

\definecolor{warnbg}{RGB}{255,243,205}
\definecolor{tipbg}{RGB}{222,240,255}

\begin{document}

{\LARGE\bfseries «Спокуха»: подготовка к вопросам жюри}

\vspace{0.6em}
Это рабочий документ для команды. Вопросы отобраны по одному принципу: жюри
видит только слайды и слышит только выступление, поэтому спрашивать будут про
то, что на слайдах написано крупно, и про то, что в бизнесе всегда спрашивают,
то есть про деньги, риски, сроки и поддержку. Сначала идут самые вероятные
вопросы, потом вопросы по отдельным слайдам, потом узкие технические вопросы
на случай, если в жюри окажется профильный специалист. В конце собраны все
числа проекта в таблицах, чтобы можно было быстро найти любую цифру прямо во
время ответа.

\vspace{0.6em}
\colorbox{warnbg}{\parbox{\dimexpr\textwidth-2\fboxsep}{%
\textbf{Три расхождения между слайдами и отчётом. Слайды менять нельзя,
поэтому снимаем голосом.}
\smallskip

\textbf{Первое, слайд 11.} На слайде написано: 1053 алерта, 36 уведомлений,
снижение шума 96,6 процента, полнота 7 из 8. В отчёте стоят другие числа:
44 уведомления, 95,8 процента и полнота 8 из 8. Слайд показывает более раннюю
версию, которая теряла один инцидент из восьми. Мы это исправили, и система
стала отправлять дежурному 44 сообщения вместо 36, то есть на восемь сообщений
больше за те же два часа. Хорошая новость в том, что слайд честно признаёт
полноту 7 из 8, поэтому наш рассказ выглядит не как оправдание, а как
продолжение работы.
\smallskip

\textbf{Второе, слайд 13.} В дереве КПЭ написано, что уведомлений на инцидент
стало 4,5, а ошибочных адресаций 10. Актуальные числа: 5,5 и 13. Причина та
же самая. Экономический итог при этом не изменился, потому что эффект считается
не от этих чисел напрямую.
\smallskip

\textbf{Третье, слайд 5.} В описании пилота написано «Zabbix + SolarWinds», а
на слайде 14 и в отчёте у нас Zabbix и Prometheus. Слайд 5 описывает контур
заказчика, каким он задан во вводных кейса, а мы на своём стенде подняли
Prometheus, потому что доступа к SolarWinds у нас не было. Для архитектуры это
безразлично: любой источник подключается отдельным адаптером.}}

\tableofcontents

\section{Самые вероятные вопросы}

\qa{Что из показанного реально работает, а что нарисовано?}
{Работает сквозной путь: событие приходит из Zabbix или Prometheus, приводится
к единой схеме, проходит корреляцию, превращается в уведомление и доставляется
в TrueConf живому человеку. Работает веб-интерфейс со сводкой и каталогом
систем. Работает локальная языковая модель. Мы специально сделали так, чтобы
границу было видно без объяснений: на всех схемах сплошная рамка это то, что
работает на стенде, а штриховая рамка это спроектированная, но ещё не
реализованная часть. В самом интерфейсе разделы, не подключённые к данным,
помечены баннером «Демонстрационные настройки». Мы принципиально нигде не
показываем правдоподобные выдуманные данные.}

\qa{Зачем разрабатывать своё, если на рынке есть готовые продукты?}
{Ни один готовый продукт не проходит по ограничениям этого кейса
одновременно. Зарубежные облачные сервисы вроде PagerDuty требуют вывода
данных наружу и сейчас в России недоступны. Встроенные средства самих систем
мониторинга видят только свои события, а нам нужна корреляция между семью
системами сразу. Российские ITSM-платформы обычно требуют реального Active
Directory и не умеют штатно работать с TrueConf. Плюс отдельное требование
заказчика: языковая модель должна работать внутри контура, а не через внешний
интерфейс. Наше решение закрывает всё это сразу.}

\qa{Вы убрали 96 процентов уведомлений. Где гарантия, что вы не выбросили
настоящую аварию?}
{Это главный вопрос к любой системе фильтрации, и мы измеряем его отдельной
метрикой. Она называется полнота и показывает, сколько инцидентов из
имеющихся система вообще заметила. На слайде честно написано семь из восьми,
то есть та версия действительно теряла один инцидент. Мы разобрались, почему
так вышло, и добавили правило: если критичность растёт выше того уровня,
который система уже видела по этому сервису, уведомление отправляется всегда.
После этого полнота стала восемь из восьми. Заплатили мы за это тем, что
система стала отправлять дежурному 44 сообщения вместо 36, то есть на восемь
сообщений больше за те же два часа, а снижение шума уменьшилось с 96,6 до
95,8 процента. Мы сознательно пошли на этот размен, потому что для системы
оповещения пропущенная авария несопоставимо дороже восьми лишних сообщений.}

\qa{Сколько нужно денег и времени, чтобы довести это до промышленной
эксплуатации?}
{Пилот уже стоил 1920 человеко-часов, это шесть человек за восемь недель, или
6,7 миллиона рублей по ставке кейса. Для промышленной эксплуатации нужно
закрыть список задач первого приоритета, он у нас составлен и приоритизирован:
аутентификация источников событий, приоритетные очереди для критичных алертов,
хранилище секретов, резервный путь оповещения на случай отказа платформы,
контроль молчания источников. По нашей оценке это ещё один сопоставимый цикл
работ. Дальше идёт сопровождение, это 30 часов в месяц, и инфраструктура,
0,76 миллиона рублей в год при аренде.}

\qa{Кто будет всё это поддерживать после вас?}
{Решение спроектировано так, чтобы повседневная эксплуатация не требовала
разработчика. Роли, правила классификации, маршрутизация и пороги это
настройки в таблицах, а не код: добавить роль или изменить правило может
администратор через интерфейс. Разработчик нужен только для подключения
принципиально нового источника, и это написание одного адаптера, ядро при этом
не меняется. Заложенная нами трудоёмкость сопровождения это 30 часов в месяц,
то есть примерно четверть ставки одного человека.}

\qa{У вас пилот на 87 пользователей и четыре филиала. Что будет, когда
масштабируете на всю Группу?}
{Мы закладывали рост с самого начала, и на слайде 14 показано, что каждый из
трёх узлов растёт по своей оси. Узел приложений растёт горизонтально, просто
добавлением копий, потому что они делят нагрузку через Kafka. Узел данных
растёт вертикально, ему нужны процессор и диск. Узел искусственного интеллекта
растёт добавлением видеокарт. Подключение нового филиала или новой системы
мониторинга не требует остановки платформы и не требует переработки ядра.
Честно назовём и ограничение: сейчас это масштабирование виртуальными
машинами, а не оркестратором, поэтому оно даёт производительность, но не даёт
отказоустойчивости. Для непрерывной работы понадобится оркестратор, и это у
нас записано в план.}

\qa{Окупаемость 0,8 года выглядит слишком красиво. Откуда такие числа?}
{Почти весь эффект даёт не экономия рабочего времени, а сокращение простоя.
Стоимость простоя не наша, она задана вводными кейса: 150 миллионов рублей в
сутки, то есть 6,25 миллиона рублей в час. При такой цене достаточно
предотвратить всего полтора часа простоя за год, чтобы получить
9,4 миллиона рублей эффекта. Именно поэтому окупаемость такая быстрая. И
обратная сторона та же самая: если мы пропустим один серьёзный инцидент,
годовой эффект обнулится. Поэтому полнота обнаружения у нас стоит в дереве
КПЭ как контрольный показатель, а не как второстепенная метрика.}

\qa{Простой установки НПЗ стоит 150 миллионов в сутки, но вы же делаете
систему оповещения. При чём тут НПЗ?}
{Мы не приписываем себе весь простой. Мы считаем только ту его часть, которая
возникает из-за задержки реакции, то есть пока авария тонет в потоке
уведомлений и пока она идёт не тому дежурному. Во вводных кейса сказано, что
простой остаётся незамеченным в среднем 14 минут. Мы закладываем сокращение
на 15 минут на инцидент и всего шесть значимых инцидентов в год. Это
осознанно скромная оценка: полтора часа простоя за целый год.}

\qa{Не опасно ли доверять оповещение искусственному интеллекту? Что если он
ошибётся?}
{У нас искусственный интеллект в принципе не может испортить оповещение,
потому что он не принимает решений. Все решения принимает обычная система
правил: что считать критичным, что дублем, кого уведомить и когда отправить.
Модель только добавляет пояснение и формулирует понятный текст, причём
работает параллельно и не задерживает доставку. Если модель ошиблась,
недоступна или вернула бессмыслицу, уведомление всё равно уйдёт, просто с
текстом, собранным из полей алерта. Это проверяется автоматическим тестом.
И отдельно: модели запрещено понижать или скрывать критичные алерты, такое
решение принимается только правилами.}

\qa{Данные никуда не уходят? Что с импортозамещением?}
{Наружу не уходит ничего, включая работу языковой модели. Модель работает на
нашем сервере, а сетевые правила запрещают ИИ-зоне выход в интернет. Внешних
облачных сервисов в решении нет вообще, лицензионных отчислений тоже нет:
Python, Kafka, TimescaleDB, nginx и vLLM это свободное программное
обеспечение, а модель Qwen открытая и работает локально. Персональные данные
мы не обрабатываем вовсе, в базе только технические идентификаторы.}

\qa{Почему TrueConf? А если мы захотим другой мессенджер?}
{TrueConf это обязательное требование кейса, поэтому он сделан как штатный,
обязательный канал доставки. Но архитектурно канал доставки это отдельный
компонент: диспетчер публикует запрос на уведомление, а адаптер канала его
забирает. Добавить электронную почту, SMS или другой мессенджер это написать
ещё один такой адаптер, ядро при этом не трогается.}

\qa{Чем вы отличаетесь от других команд, которые делают тот же кейс?}
{Скорее всего тем, что мы всё измерили и показываем неудобные числа тоже.
Мы не заявляем снижение шума, а измеряем его на воспроизводимом наборе
данных и вместе с ним показываем полноту обнаружения и точность адресации,
включая то, что 13 уведомлений из 44 пока уходят не тому дежурному. Мы честно
отмечаем на схемах, что реализовано, а что только спроектировано, и в
интерфейсе помечаем макеты баннером. И у нас работающий сквозной путь на
реальном стенде, а не презентация.}

\section{Вопросы по конкретным слайдам}

\subsection{Слайд 2, команда}

\qa{Кто в команде чем занимался?}
{Шесть человек: архитектор отвечал за структуру решения и границы контуров,
три разработчика делали ядро обработки, интерфейс и интеграции, специалист по
информационной безопасности вёл модель угроз и реестр рисков, аналитик
занимался требованиями, матрицей трассировки и экономикой. Такой состав дал
нам возможность смотреть на задачу не только с технической стороны, но и со
стороны пользователя и реальной эксплуатации.}

\subsection{Слайд 5, предпосылки и бизнес-требования}

\qa{Откуда взялись 5400 событий, 65 процентов шума и 3,2 миллиона за
инцидент?}
{Это вводные данные кейса, их дал заказчик, мы их не придумывали. Наши
собственные измерения это отдельная история, они сделаны на нашем наборе
данных, и там другие абсолютные числа, потому что это наш пилотный периметр,
а не все системы заказчика.}

\qa{На слайде написано Zabbix и SolarWinds, а показываете вы Prometheus.}
{На слайде описан контур заказчика, каким он задан во вводных кейса. Доступа
к SolarWinds у нас не было, поэтому на стенде мы подняли Prometheus вместе с
Alertmanager, это распространённая система мониторинга с другим форматом
сообщений. Для нас важно было проверить именно разноформатность, то есть что
два непохожих источника сводятся к единой схеме. SolarWinds подключается точно
так же, отдельным адаптером, ядро при этом не меняется.}

\subsection{Слайд 7, архитектура}

\qa{Почему Active Directory стоит сбоку и не участвует в аутентификации?}
{Это ограничение кейса: доступа к домену у нас нет. Мы превратили ограничение
в архитектурное решение: вход выполняется через TrueConf по протоколу
OAuth 2, тот самый TrueConf, который у нас и так обязателен как канал
доставки. Active Directory остался справочником подразделений и команд, причём
платформа в него только читает и никогда ничего не записывает. Когда домен
станет доступен, подключить его можно будет без переделки, потому что вход у
нас отдельный компонент.}

\subsection{Слайд 9, обработка события}

\qa{Что такое «детерминированное ядро» простыми словами и почему это
преимущество?}
{Это означает, что решения принимает система понятных правил, а не нейросеть.
Правило можно прочитать, проверить и поправить, и оно всегда даёт один и тот
же результат на одних и тех же данных. Практическая польза простая: когда
система ошиблась, инженер открывает решение, видит, какое именно правило
сработало, и точечно его донастраивает. С нейросетью в этой роли пришлось бы
гадать. Плюс правила работают всегда, даже когда модель недоступна.}

\subsection{Слайд 11, искусственный интеллект}

\qa{На слайде полнота 7 из 8. Получается, вы потеряли инцидент?}
{Та версия, которая на слайде, действительно теряла один инцидент из восьми, и
мы это не прячем. Причина оказалась понятной: второй инцидент начинался, когда
первый ещё расходился каскадом по тем же сервисам, и система считала его
продолжением первого. Мы добавили правило эскалации: если критичность по сервису
выросла выше того уровня, который система уже видела, уведомление отправляется
обязательно, даже если до этого система решила придержать сигнал. Сейчас
полнота восемь из восьми. Стоило это восьми дополнительных сообщений дежурному
за два часа, было 36, стало 44. Половина из этих восьми попала правильному
адресату, так что размен нас полностью устраивает.}

\qa{Зачем вам вообще искусственный интеллект, если решения принимают правила?}
{Он делает уведомление понятным человеку. Правила определяют, что случилось и
кого звать, а модель добавляет к этому связное объяснение, гипотезу о
первопричине и первый шаг проверки. Дежурный получает не строку с метрикой и
порогом, а нормальный текст. При этом мы сознательно оставили модель в роли
помощника: если убрать её совсем, система продолжит работать, просто
уведомления станут суше.}

\subsection{Слайд 12, информационная безопасность}

\qa{Вы сами пишете, что webhook не аутентифицирован. Разве такое можно
внедрять?}
{В таком виде нельзя, и именно поэтому мы вынесли это первым пунктом в
остаточные риски. На изолированном стенде это допустимо, в промышленной
эксплуатации недопустимо. Мы предпочли показать честный список того, что
осталось закрыть, а не рисовать зелёные галочки. Всего у нас 23 сценария
угроз: шесть закрыты полностью, одиннадцать частично, три требуют работы до
внедрения, три неактуальны.}

\qa{Почему именно угрозы искусственного интеллекта закрыты полностью?}
{Потому что мы не дали модели никаких полномочий. Она не пишет в базу, не
меняет настройки, не отправляет уведомления и не может понизить критичность.
Текст алерта приходит к ней как данные, отдельно от инструкции, а ответ
принимается только в пределах заранее заданного списка значений. Это не
навешенное сверху средство защиты, это следствие архитектуры, поэтому такие
угрозы закрыты не частично, а полностью.}

\subsection{Слайд 13, дерево КПЭ}

\qa{Как вы вообще перешли от уведомлений к рублям?}
{В два шага, и оба видны на слайде. Сначала измеренные метрики превращаются в
физические величины: сэкономленные человеко-часы и сокращённые часы простоя.
Потом эти величины умножаются на обязательные ставки кейса, то есть на
2000 рублей за час работы сотрудника и 6,25 миллиона рублей за час простоя.
Получается 2,45 миллиона рублей на трудозатратах и 9,38 миллиона на простое,
вместе 11,8 миллиона рублей в год.}

\qa{Ваша экономия рабочего времени выглядит скромно на фоне 48 человеко-часов
в сутки из вводных.}
{Так и есть, и это сделано намеренно. Наивный расчёт дал бы огромное число,
но он был бы неправдой: сегодня эти часы на самом деле не тратятся, поток
уведомлений просто игнорируется, и именно поэтому теряются аварии. Поэтому мы
ограничили эффект реальным фондом времени дежурных и взяли только ту часть,
которая действительно высвобождается. Получилось скромно, зато защитимо.}

\subsection{Слайд 14, инфраструктура}

\qa{Зачем вам отдельный сервер с видеокартой? Это дорого.}
{Видеокарта нужна потому, что модель постоянно держит в видеопамяти около
14 гигабайт, независимо от того, есть нагрузка или нет. Держать такой ресурс
внутри общего сервера расточительно, поэтому мы вынесли его отдельно. По
деньгам это 45 тысяч рублей в месяц из 63 тысяч, то есть 71 процент стоимости
всей инфраструктуры. И тут есть важное следствие: если заказчик решит, что
искусственный интеллект ему не нужен, этот узел просто выключается, система
продолжает работать, а инфраструктура дешевеет с 0,76 до 0,22 миллиона рублей
в год.}

\qa{Три виртуальные машины на 87 пользователей. Не избыточно ли?}
{Нет, потому что размер узлов определяется не числом пользователей, а
профилем нагрузки. Мы сняли замер работающего стенда: наши собственные
приложения занимают меньше 150 мегабайт памяти, а основные потребители это
шина сообщений и модель. Разделение на три узла нужно не ради запаса, а ради
того, чтобы эти три разные нагрузки не мешали друг другу и росли независимо.
Запас по узлу приложений при этом действительно большой, примерно
пятидесятикратный, и он бесплатный.}

\section{Узкие технические вопросы}

Этот раздел на случай, если в жюри окажется профильный специалист.

\qa{Почему Kafka, а не более простая очередь или прямые вызовы?}
{Нам нужна была журналируемая шина, где сообщения хранятся, а не исчезают
после доставки. Это даёт три вещи: недоступность потребителя не теряет
события, обработчики масштабируются горизонтально за счёт групп потребителей,
и приём событий отделён от их обработки, поэтому всплеск нагрузки переживается
буфером. Прямые синхронные вызовы создавали бы обратное давление на приём, а
это запрещено требованиями.}

\qa{Почему TimescaleDB?}
{События и алерты это временные ряды, и гипертаблицы сразу дают сжатие и
политики хранения: у нас события и алерты живут 90 дней, а решения
искусственного интеллекта 400 дней. При этом под капотом остаётся обычный
PostgreSQL, поэтому реляционная часть с ролями и связями живёт в той же базе.
В итоге одна СУБД вместо двух.}

\qa{Как именно работает корреляция?}
{Каждый алерт получает ключ: явный идентификатор корреляции, если он есть,
иначе сервис, иначе пара из источника и метрики. На каждый ключ в пределах
временного окна работает конечный автомат, который решает, что это: повтор,
дополнительное свидетельство той же аварии, нестабильный источник или новая
авария, требующая уведомления. Важная деталь: состояние не хранится в
изменяемом поле, оно каждый раз выводится проигрыванием всего окна заново,
поэтому оно не может разойтись с историей.}

\qa{Что происходит при отказе базы данных?}
{Обработка и доставка продолжаются, сбой записи не выходит в поток обработки,
и это подтверждено тестом. Теряется только история за время отказа. Мы
считаем это правильным приоритетом для системы оповещения, но не считаем
идеалом: подтверждение приёма только после надёжного сохранения стоит у нас в
списке первого приоритета.}

\qa{Как вы измеряли качество и почему это не подгонка?}
{Мы построили набор из 1053 алертов, где 545 значимых и 508 шумовых, восемь
инцидентов и два часа модельного времени с фиксированным начальным значением
генератора. Разметка не проставлена руками, она выводится теми же правилами,
которые порождают события. Ключевой момент честности: разметка не попадает в
конвейер. Корреляция не видит идентификатор инцидента, иначе она сгруппировала
бы алерты именно по нему и показала бы идеальный результат по совершенно
неверной причине. Границу измерения мы тоже называем прямо: один и тот же
источник задаёт и шум, и критерий шума, поэтому это честный стенд для
сравнения версий между собой, но не доказательство эффективности на реальном
трафике.}

\qa{Что с тестированием?}
{145 тестов платформы проходят успешно, один пропущен, потому что требует
развёрнутой модели, четыре помечены как известные ограничения. Отдельно
покрыты диспетчер доставки и генератор событий. Тестами зафиксированы именно
те свойства, на которых держится экономика: отключение модели не останавливает
доставку, ответ вне словаря не применяется, отказ базы не прерывает обработку,
некорректный webhook отклоняется.}

\qa{Точность 31,8 процента, это же мало.}
{Мало, и мы измерили это сами, а не ждали, пока заметят. Причина конкретная:
13 уведомлений из 44 уходят дежурному того сервиса, который сигнализирует о
проблеме, а не того, который сломался. Лечится это учётом первопричины и
данными сервисного каталога при маршрутизации, и это основная задача
следующего этапа. Цель 60 процентов, и проверяться она будет на том же наборе
данных.}

\qa{Как подключить новую систему мониторинга?}
{Написать адаптер, ядро не трогается. Есть два способа подключения: либо
источник сам вызывает нас по webhook, либо мы опрашиваем его по расписанию.
Формат данных полностью принадлежит адаптеру. Внутренняя схема события
минимальная и открытая: обязательны только источник и статус, а остальные поля
вендора сохраняются как есть, поэтому миграция базы не нужна.}

\section{Если вопрос застал врасплох}

\colorbox{tipbg}{\parbox{\dimexpr\textwidth-2\fboxsep}{%
Три правила на такой случай. Первое: если числа нет в голове, не выдумывайте,
скажите, что оно есть в отчёте, и назовите порядок величины. Второе: если
вопрос про то, чего мы не сделали, отвечайте прямо, что не сделали, и сразу
называйте, в какой очереди это стоит и почему именно там. Слабость,
названная первой, перестаёт быть слабостью. Третье: если вопрос выходит за
рамки того, что мы проверяли, так и скажите, и добавьте, каким измерением это
можно было бы проверить.}}

\section{Словарик: что означают наши термины}

Здесь простыми словами объяснено то, что звучит в выступлении и в ответах.
Пригодится, если вопрос задаст нетехнический член жюри и объяснять придётся
на ходу.

\begin{longtable}{|L{3.6cm}|L{12.9cm}|}
\hline
\textbf{Термин} & \textbf{Что это значит простыми словами} \\ \hline
Событие & Любая запись, которую прислала система мониторинга. Само по себе оно ещё ничего не значит \\ \hline
Алерт & Событие, которое система мониторинга сочла достойным внимания: превышен порог, упал сервис, кончилось место \\ \hline
Инцидент & Одна настоящая авария. Одна авария обычно порождает десятки алертов из разных систем, поэтому инцидентов всегда сильно меньше \\ \hline
Уведомление & Сообщение, которое реально уходит живому дежурному. Наша задача в том, чтобы из тысячи алертов получилось несколько уведомлений \\ \hline
Шум & Алерты, которые не требуют реакции человека: повторы, мелочь, следствия уже известной аварии \\ \hline
\textbf{Полнота} & Про сколько настоящих аварий система вообще сообщила. У нас 8 из 8: ни одна авария не осталась незамеченной. Это защита от самой дорогой ошибки, то есть от молчания при настоящей аварии \\ \hline
\textbf{Точность} & Какая доля отправленных сообщений оказалась полезной. У нас 14 верных из 44 отправленных, это 31,8 процента. Точность и полнота тянут в разные стороны: отправляя всё подряд, получишь идеальную полноту и ужасную точность \\ \hline
Дедупликация & Отсечение повторов: одна и та же проблема повторяется каждую минуту, а человеку сообщаем один раз \\ \hline
Корреляция & Склейка разных сигналов, у которых одна причина. Упал сервер, и об этом закричали и Zabbix, и Prometheus, и сервисы сверху: это один инцидент, а не пять \\ \hline
Эскалация & Правило, по которому мы обязаны сообщить, даже если до этого решили промолчать. У нас оно срабатывает, когда критичность выросла выше уже виденного уровня \\ \hline
Детерминированный & Работающий по понятным правилам, а не по догадке нейросети. Правило можно прочитать, проверить и поправить, и оно всегда даёт один и тот же результат \\ \hline
Fallback, резервный путь & Что система делает, когда что-то сломалось. У нас: если модель недоступна, текст уведомления собирается из полей алерта, и уведомление всё равно уходит \\ \hline
Маршрутизация & Определение, кому именно отправить: по сервису находим команду, по команде дежурного, по дежурному канал доставки \\ \hline
FTE & Условный сотрудник на полной ставке. Ставка кейса 2000 рублей за час его работы \\ \hline
NPV, чистая приведённая стоимость & Сколько проект приносит за три года с поправкой на то, что рубль завтра дешевле рубля сегодня. У нас 13,9 миллиона \\ \hline
DPP, срок окупаемости & Через сколько проект вернёт вложенное. У нас 0,8 года \\ \hline
PI, индекс доходности & Во сколько раз отдача больше вложений. У нас 3,1 \\ \hline
\end{longtable}

\section{Все числа проекта в таблицах}

\subsection{Вводные данные кейса, описывающие проблему заказчика}
\begin{longtable}{|L{9cm}|L{7.5cm}|}
\hline
\textbf{Параметр} & \textbf{Значение} \\ \hline
Систем мониторинга & 7 \\ \hline
Событий в сутки & около 5400 \\ \hline
Доля шума & 65 процентов \\ \hline
Ручной разбор & около 48 человеко-часов в сутки \\ \hline
Простой остаётся незамеченным в среднем & 14 минут \\ \hline
Стоимость одного инцидента & около 3,2 миллиона рублей \\ \hline
Охват пилота & 4 филиала, около 87 пользователей \\ \hline
Источники пилота по вводным & Zabbix и SolarWinds; на нашем стенде Zabbix и Prometheus \\ \hline
Требуемое время доставки уведомления & 2--3 минуты \\ \hline
Обновление дашборда & не дольше 5 минут; новое событие на панели до 15 секунд \\ \hline
\end{longtable}

\subsection{Обязательные ставки для финансовой модели}
\begin{longtable}{|L{9cm}|L{7.5cm}|}
\hline
\textbf{Параметр} & \textbf{Значение} \\ \hline
Стоимость FTE, применяется при расчёте эффектов & 2000 рублей в час \\ \hline
Час работы проектной команды, применяется при расчёте затрат & 3500 рублей в час \\ \hline
Простой установки НПЗ & 150 миллионов рублей в сутки, то есть 6,25 миллиона в час \\ \hline
\end{longtable}

\subsection{Качество обработки: что на слайде и что актуально}
\begin{longtable}{|L{7cm}|L{9.5cm}|}
\hline
\textbf{Параметр} & \textbf{Значение} \\ \hline
Состав набора данных & 1053 алерта, из них 545 значимых и 508 шумовых; 8 инцидентов; 2 часа модельного времени \\ \hline
Без фильтрации & 1053 уведомления; снижение шума 0 процентов; полнота 8 из 8; точность 20,5 процента \\ \hline
Версия на слайде 11 & 36 уведомлений; снижение шума 96,6 процента; полнота \textbf{7 из 8}; точность 27,8 процента \\ \hline
Актуальная версия & \textbf{44 уведомления; снижение шума 95,8 процента; полнота 8 из 8; точность 31,8 процента} \\ \hline
Как распределились 44 уведомления & 14 адресованы верно, 13 ушли не тому дежурному, 17 остались шумовыми, причём из них лишь одно относится к окну обслуживания \\ \hline
Как распределилась базовая линия & 216 верных, 329 не тому адресату, 508 шумовых \\ \hline
Что дало правило эскалации & Полнота выросла с 7 из 8 до 8 из 8. Дежурный стал получать на 8 сообщений больше за два часа, было 36, стало 44; из этих восьми четыре попали верному адресату. Снижение шума уменьшилось с 96,6 до 95,8 процента, а точность выросла с 27,8 до 31,8 процента \\ \hline
\end{longtable}

\subsection{Продуктовые метрики: база, достигнуто, цель}
\begin{longtable}{|L{5.6cm}|L{2.6cm}|L{2.9cm}|L{2.4cm}|}
\hline
\textbf{Метрика} & \textbf{База} & \textbf{Достигнуто} & \textbf{Цель} \\ \hline
Уведомлений за 2 часа набора данных & 1053 & 44 & не более 50 \\ \hline
Уведомлений в час & 527 & 22 & не более 25 \\ \hline
Уведомлений на один инцидент & 132 & 5,5 (на слайде 4,5) & не более 6 \\ \hline
Шумовых уведомлений & 508 & 17 & не более 10 \\ \hline
Уведомлений не тому адресату & 329 & 13 (на слайде 10) & не более 5 \\ \hline
Доля полезных уведомлений & 20,5 \% & 31,8 \% & не менее 60 \% \\ \hline
Полнота обнаружения инцидентов & 8 из 8 & 8 из 8 & 8 из 8 \\ \hline
Время до первого уведомления & не измерялось & не измерялось & не более 15 с \\ \hline
Доля немаршрутизируемых алертов & не измерялось & не измерялось & не более 5 \% \\ \hline
Доставка при недоступности ИИ & не измерялось & 100 \% & 100 \% \\ \hline
\end{longtable}

\subsection{Собственные допущения финансовой модели}
\begin{longtable}{|L{6.4cm}|L{3cm}|L{6.8cm}|}
\hline
\textbf{Параметр} & \textbf{Значение} & \textbf{Обоснование} \\ \hline
Поток алертов пилотного периметра & 12 636 штук в сутки & Измеренный набор данных, умноженный на 12 \\ \hline
Среднее время реакции на уведомление & 2 минуты & Время открыть, прочитать и квалифицировать оповещение \\ \hline
Дежурные смены & 2 смены по 8 часов & Типовая организация дежурства \\ \hline
Доля времени дежурного на разбор оповещений & 30 процентов & Ограничивает эффект фактическим фондом времени \\ \hline
Доля высвобождаемого времени & 70 процентов & Часть разбора остаётся при любом снижении шума \\ \hline
Значимых инцидентов с простоем в год & 6 & Оценка пилотного периметра, уточняется по журналу аварий \\ \hline
Сокращение времени обнаружения на инцидент & 15 минут & Следствие снижения числа уведомлений и адресной доставки \\ \hline
Состав проектной команды & 6 человек & Архитектор, три разработчика, специалист по ИБ, аналитик \\ \hline
Трудоёмкость пилота & 1920 часов & 6 человек, 8 недель, по 40 часов \\ \hline
Трудоёмкость сопровождения & 30 часов в месяц & Поддержка адаптеров, правил и конфигурации \\ \hline
Ставка дисконтирования & 20 процентов & Типовая для проектов внутренней ИТ-автоматизации \\ \hline
Горизонт расчёта & 3 года & Срок жизни пилотного решения до пересмотра \\ \hline
\end{longtable}

\subsection{Расчёт эффектов}
\begin{longtable}{|L{10.4cm}|L{6cm}|}
\hline
\textbf{Расчёт} & \textbf{Результат} \\ \hline
Трудозатраты: 16 человеко-часов в сутки, из них 30 процентов даёт 4,8; из них высвобождается 70 процентов, это 3,36 в сутки; за год примерно 1226 человеко-часов; умножаем на 2000 рублей & \textbf{2,45 миллиона рублей в год} \\ \hline
Простой: 6 инцидентов по 0,25 часа дают 1,5 часа в год; умножаем на 6,25 миллиона рублей за час & \textbf{9,38 миллиона рублей в год} \\ \hline
Суммарный операционный эффект & \textbf{11,8 миллиона рублей в год} \\ \hline
\end{longtable}

\subsection{Затраты}
\begin{longtable}{|L{6.4cm}|L{4.4cm}|L{5.2cm}|}
\hline
\textbf{Статья} & \textbf{Расчёт} & \textbf{Сумма} \\ \hline
Разработка и внедрение пилота & 1920 часов по 3500 рублей & 6,7 миллиона рублей, разово \\ \hline
Сопровождение и развитие & 360 часов в год по 3500 рублей & 1,3 миллиона рублей в год \\ \hline
Инфраструктура, аренда трёх узлов & 63 тысячи рублей на 12 месяцев & 0,76 миллиона рублей в год \\ \hline
Лицензии внешних сервисов & отсутствуют & 0 \\ \hline
\textbf{Итого} & & \textbf{6,7 миллиона рублей разово и 2,0 миллиона рублей в год} \\ \hline
\end{longtable}

Инфраструктура по месяцам: узел приложений (4 vCPU, 8 ГБ памяти, 50 ГБ диска)
стоит 6 тысяч рублей, узел данных и шины (8 vCPU, 16 ГБ памяти, 200 ГБ SSD)
стоит 12 тысяч рублей, узел искусственного интеллекта с видеокартой на 16 ГБ
стоит 45 тысяч рублей, всего 63 тысячи рублей в месяц. На долю ИИ-узла
приходится 71 процент. Если функции искусственного интеллекта выключить,
инфраструктура обойдётся в 0,22 миллиона рублей в год.

\subsection{Инвестиционные показатели при аренде, миллионы рублей}
\begin{longtable}{|L{6.2cm}|c|c|c|c|}
\hline
\textbf{Показатель} & \textbf{Год 0} & \textbf{Год 1} & \textbf{Год 2} & \textbf{Год 3} \\ \hline
Эффект & нет & 11,8 & 11,8 & 11,8 \\ \hline
Операционные затраты & нет & 2,0 & 2,0 & 2,0 \\ \hline
Инвестиции & 6,7 & нет & нет & нет \\ \hline
Чистый поток & $-6{,}7$ & 9,8 & 9,8 & 9,8 \\ \hline
Коэффициент дисконтирования при ставке 20 \% & 1,00 & 0,83 & 0,69 & 0,58 \\ \hline
Дисконтированный поток & $-6{,}7$ & 8,2 & 6,8 & 5,7 \\ \hline
Накопленный дисконтированный поток & $-6{,}7$ & 1,5 & 8,3 & \textbf{13,9} \\ \hline
\end{longtable}

Итог: чистая приведённая стоимость \textbf{13,9 миллиона рублей}, приведённая
стоимость 20,7 миллиона, приведённые инвестиции 6,7 миллиона, индекс
доходности \textbf{3,1}, срок окупаемости \textbf{около 0,8 года}, то есть
примерно десять месяцев.

\subsection{Сравнение аренды и покупки оборудования}
\begin{longtable}{|L{7.6cm}|L{4cm}|L{4.4cm}|}
\hline
\textbf{Показатель} & \textbf{Аренда} & \textbf{Покупка} \\ \hline
Разовые инвестиции, миллионы рублей & 6,7 & 9,06 \\ \hline
Стоимость оборудования, миллионы рублей & не требуется & 2,36, из них узел с видеокартой 1,66 \\ \hline
Периодические затраты, миллионы рублей в год & 2,0 & 1,65 \\ \hline
Чистая приведённая стоимость за 3 года & 13,9 & 12,3 \\ \hline
Срок окупаемости, лет & 0,8 & 1,1 \\ \hline
\end{longtable}

\subsection{Чувствительность модели}
\begin{longtable}{|L{9.6cm}|L{6.9cm}|}
\hline
\textbf{Утверждение} & \textbf{Число} \\ \hline
Сколько простоя нужно предотвратить, чтобы покрыть все затраты первого года, а это 8,7 миллиона рублей & 1,4 часа в год \\ \hline
Чистая приведённая стоимость, если полностью обнулить эффект от экономии трудозатрат & около 8,8 миллиона рублей \\ \hline
Цена ошибки в один час простоя & 6,25 миллиона рублей в год \\ \hline
Последствие одного пропущенного бизнес-критичного инцидента & может обнулить весь годовой эффект \\ \hline
\end{longtable}

\subsection{Замер потребления ресурсов на пилотном стенде}
Стенд это виртуальная машина с 16 vCPU, 62 ГБ оперативной памяти, 133 ГБ
диска и видеокартой Tesla T4 на 15 ГБ. На ней работает 15 контейнеров, аптайм
на момент замера пять суток. Замер снят в режиме ожидания, поэтому показывает
базовое, а не пиковое потребление.
\begin{longtable}{|L{5.6cm}|L{2.2cm}|L{2.4cm}|L{5.6cm}|}
\hline
\textbf{Компонент} & \textbf{Процессор} & \textbf{Память} & \textbf{Примечание} \\ \hline
Платформа, основной процесс & 0,4 \% & 65 МБ & Приём, корреляция, программный интерфейс \\ \hline
Веб-интерфейс & менее 0,1 \% & 13 МБ & Статика и обратный прокси \\ \hline
TimescaleDB & 0,8 \% & 384 МБ & Растёт вместе с объёмом истории \\ \hline
Kafka & 181 \% & 1,2 ГБ & В простое занимает около 1,8 ядра \\ \hline
Генератор событий & 0,4 \% & 40 МБ & Существует только на стенде \\ \hline
Инференс vLLM & 1,2 \% & 4,0 ГБ & Плюс 13,6 из 15,4 ГБ видеопамяти \\ \hline
TrueConf Server & 0,5 \% & 1,25 ГБ & Внешняя система, обычно уже есть \\ \hline
Контроллер домена & 0,2 \% & 255 МБ & Синтетический справочник стенда \\ \hline
Zabbix, четыре контейнера & 1,1 \% & 945 МБ & Источник событий, обычно уже есть \\ \hline
Prometheus & менее 0,1 \% & 44 МБ & Источник событий и метрик \\ \hline
Сервис CMDB & 0,1 \% & 33 МБ & Развёрнут, но пока не используется \\ \hline
TLS-фронт & менее 0,1 \% & 16 МБ & Терминация TLS \\ \hline
\textbf{Итого} & \textbf{около 1,9 ядра} & \textbf{около 9 ГБ} & Диск: 38 ГБ образов и 2,2 ГБ томов \\ \hline
\end{longtable}

\subsection{Целевая конфигурация узлов}
\begin{longtable}{|L{3.2cm}|L{5.4cm}|L{7.4cm}|}
\hline
\textbf{Узел} & \textbf{Конфигурация} & \textbf{Чем ограничен и как растёт} \\ \hline
Приложения & 4 vCPU, 8 ГБ, 50 ГБ & Ограничен числом экземпляров, а не ресурсами; растёт горизонтально; запас примерно пятидесятикратный \\ \hline
Данные и шина & 8 vCPU, 16 ГБ, 200 ГБ SSD & Ограничен процессором и диском; растёт вертикально и вместе со сроком хранения \\ \hline
Искусственный интеллект & 8 vCPU, 32 ГБ, видеокарта 16 ГБ & Ограничен видеопамятью и числом одновременных запросов; растёт добавлением узлов с видеокартами \\ \hline
Интеграции & Ресурсы заказчика & Zabbix, Prometheus, каталог и TrueConf обычно уже развёрнуты, в расчёт затрат не входят \\ \hline
\end{longtable}

\subsection{Сводные числа по информационной безопасности}
\begin{longtable}{|L{7cm}|L{9.5cm}|}
\hline
\textbf{Параметр} & \textbf{Значение} \\ \hline
Разобрано сценариев угроз & 23 \\ \hline
Нейтрализовано полностью & 6, включая все угрозы ИИ-контура \\ \hline
Нейтрализовано частично & 11 \\ \hline
Требует работы до внедрения & 3: неаутентифицированный webhook У-01, подмена образа или весов модели У-16, отказ единственного узла У-18; рядом стоит контроль молчания источника У-19 \\ \hline
Неактуальны для нашей архитектуры & 3: персональные данные, база знаний, обращения к CMDB и каталогу \\ \hline
Категорий нарушителей & 7, от внешнего без доступа к сети до автора текста алерта \\ \hline
Границ доверия & 6; граница между ядром и ИИ реализована полностью \\ \hline
Тесты платформы & 145 пройдено, 1 пропущен, 4 отмечены как известные ограничения \\ \hline
Утверждений о защите, подтверждённых автотестами & 7 \\ \hline
Сроки хранения & события и алерты 90 дней, решения ИИ 400 дней \\ \hline
\end{longtable}

\subsection{Прочие технические числа}
\begin{longtable}{|L{7cm}|L{9.5cm}|}
\hline
\textbf{Параметр} & \textbf{Значение} \\ \hline
Время вызова модели для формулировки уведомления & около 3,4 секунды на видеокарте Tesla T4 \\ \hline
Словарь приоритетов, доступный модели & p1, p2, p3, p4 \\ \hline
Архитектурных решений в матрице трассировки & 17 \\ \hline
Требований к интеллектуальной обработке & 14 позиций: 9 реализовано, 4 частично, 1 отнесена к целевой архитектуре \\ \hline
Этапов плана внедрения & 5: наблюдение, затем корреляция и маршрутизация, затем теневой и подсказывающий режимы, затем расширение периметра, затем ограниченная автоматизация \\ \hline
Ответ на некорректную нагрузку webhook & код 422 \\ \hline
Ответ нереализованных методов интерфейса & код 501, а не пустой список \\ \hline
Порты пилотного стенда & Zabbix 8080; Prometheus 9090 и 9093; генератор событий 8090; платформа 8100; веб-интерфейс 8091; уведомления 8110; бот TrueConf 8120; TrueConf Server 80, 443, 4307; TLS-фронт 8443 и 8543; база данных портов не публикует \\ \hline
\end{longtable}

\end{document}
